\begin{afterword}
\summaryhead

People talk as if Einstein had superpowers. The author thought so too, until Julian Barbour's \textsc{The End of Time} showed Einstein's work as understandable, and the author imagines Barbour quietly saying that what Einstein did isn't magic. (A 2013 edit says Barbour never said it.) John Baez's Crackpot Index scores favorable comparisons with Einstein. The author argues that the crackpot, who claims Einstein's magic, and the critic shocked by a young physicist's hope to do better than Einstein share one error: they confuse social status with research potential, and imagine Einstein's potential as rare as his fame.

Not everyone can be Einstein, the author grants. But many bright people end up ``just another Jewish genius,'' and what separates them from Einstein, the author suspects, is not brilliance but choosing an important problem, finding a worthwhile angle of attack, and persisting for years without support. Before his papers, Einstein would have shown no aura of destiny. The author, who has such a problem, meets people who ask to see that aura. Barbour saw through Einstein, and one acquires such powers by seeing that they are normal.

\responsehead

The post's account of great work is plausible, and it is not new. That great work takes an important problem, a way to attack it and years of persistence, more than raw brains, was the theme of Richard Hamming's well-known talk ``You and Your Research'' (1986). Einstein's early record fits the post's point about the missing aura: two years without a teaching post, then a junior examiner's job at the patent office. The post would have been stronger with either.

What it offers instead is thinner. The ``widespread tendency'' is shown by two items from a joke list, which show one physicist's experience rather than ``what society generally thinks.'' The Barbour quotation was always, by the post's own account, the author's reading between the lines, and the guess that people sneered at Barbour is offered with ``I kinda suspect'' and no evidence. The central claim, that what separates bright people from Einstein is problem choice rather than brilliance, is offered with ``I suspect.'' The post argues for it; it does not test it.

The post states its stake openly: the author has an important problem, an angle, years of work, and meets people who say, ``Let's see your aura of destiny, buddy.'' That line is the author's paraphrase. Read as a request for evidence of ability, it is the ``ordinarily evaluable claim'' the post itself asked for, and an ordinary evaluation asks for a record. The one feat the post offers, the AI-Box Experiment, is described on its own page as ``strictly anecdotal evidence,'' with only the outcomes revealed.

In short: a plausible account of great work, close to Hamming's and borne out by Einstein's early career, argued here from a joke list, an imagined quotation and the author's own encounters.
\end{afterword}
